\chapter{Server-Class Cluster Vision}
\label{chap:cluster-vision}

CharlotteOS explores an operating-system boundary at which the cluster is the
primary administrative and computational unit. An operator assigns a signed
workload to the cluster, declares its resources, authority, multiplicity, and
placement constraints, and leaves physical placement to cluster policy. Nodes
supply processors, memory, devices, and failure domains; replacing one should
not require rebuilding a machine-specific application installation.

Three kinds of evidence appear in this chapter and should not be confused.
The generic Raft core exercises quorum rules with configurations including
five voters, and the safety models explore three-node elections and membership
transitions. Operational QEMU validation ranges from focused two-guest
scenarios to a three-guest cluster with discovery-driven admission,
joint-consensus membership, replica placement, and cluster-wide TCP ingress.
The number of guests in a fixture is a validation choice, not a cluster-size
limit or architectural objective.

The implemented system now includes replicated naming and deployment state,
signed ELF artifacts, central-S3 artifact pickup, per-node deployment agents,
remote invocation, deterministic replica placement, basic affinity and
anti-affinity, membership and drain reconciliation, exact-generation
readiness, and a DSR-backed service VIP. Capacity-aware scheduling, stateful
cross-node migration, rolling replacement, and a production administration
plane define the next engineering boundary.

This chapter first develops the cluster-level model and then uses Linux
containers and Kubernetes/OpenShift as a counterpoint. The comparison separates
the intrinsic cost of distributed computation from machinery needed to adapt
general-purpose, machine-centred operating systems to isolated, reproducible,
movable workloads. CharlotteOS is a research prototype for testing that
separation, not a production container-platform replacement.

\section{The cluster is the computer}

\subsection{Nodes are resources, not deployment targets}

The central proposition is simple:

\begin{quote}
    The Charlotte cluster is the administrative unit; nodes contribute
    replaceable execution and storage resources.
\end{quote}

A node boots, discovers or is directed to a cluster, establishes its
identity, joins shared state, advertises resources, and receives work. Its
useful description consists of properties such as processor architecture, memory,
accelerators, attached devices, trust domain, and failure domain. If a node
fails, the operator should replace the resource and let the cluster
reconstruct the intended placement. If a workload moves, clients should
continue to address its cluster name.

That indirection now has an external TCP counterpart. Node addresses serve
protocols that intentionally target one machine. A cluster service VIP instead
addresses functionality placed somewhere in Charlotte. Direct Server Return
both distributes load and hides ingress and execution nodes behind a stable
cluster network identity. The
implemented data path and its current placement-policy boundary are described
in \autoref{sec:distributed-l2-ingress}.

Consensus membership, failure detection, routing, and ownership fencing still
use stable node identities. Those identities remain mechanisms below the
operator's application-deployment model. Interchangeability also permits
heterogeneous hardware: schedulable properties and failure-domain constraints
express relevant differences without turning nodes into hand-maintained
machine personalities.

The same idea already appears one level down in CharlotteOS. The
shard-per-core model of \autoref{chap:sitas-executor} treats cores as executors
over explicitly owned state, while the scheduler in
\autoref{chap:scheduling} moves work between them. The
cluster vision scales this relationship outward: a cluster node is analogous
to a core, cluster-visible storage is analogous to owned persistent state,
and the cluster becomes the unit within which execution is placed.

\subsection{Workloads are signed execution objects}

In the target model, cluster software is a signed execution object with a
logical identity and a deliberately small execution environment. Its signed
metadata can bind, among other things, its name, class, release, rollback
counter, permission to run concurrent instances, and a digest linking to
retained provenance.

Signature validity establishes the artifact identity and accepted origin.
Capability delegation independently determines what a running workload may
do. Each workload executes in its own address space and receives a deliberate
set of service endpoints, storage objects, devices, or administration
capabilities.

CharlotteOS therefore combines an immutable, name-bound identity at the
loading boundary with explicit authority at the execution boundary. Signing
without confinement accepts an over-powerful program; confinement without
authentication accepts code of unknown origin.

This boundary also changes what installing software means. Build systems
resolve third-party code and produce a self-contained, blessed artifact with
its evidence. Nodes avoid installation scripts and machine-local dependency
state. Reproducible builds, vulnerability policy, SBOMs, source attestation,
and signing-key custody remain supply-chain responsibilities.

\subsection{Assignment is declarative and cluster-wide}

The desired operation makes version 71 of a claims engine active in the cluster with a
desired number of instances and declared resource, authority, affinity,
anti-affinity, and failure-domain constraints.

That declaration becomes replicated desired state. A cluster controller
compares it with observed state and chooses assignments. An assigned node
obtains the immutable artifact, verifies its identity and signature, creates
an isolated execution domain, and supplies its delegated capabilities.
Rollback is another desired-state change to an earlier accepted artifact,
not a sequence of machine-repair commands.

Placement maps a component to one or more eligible nodes. A component explicitly
blessed as safe for parallel execution may have several instances. Closely
interacting components may be co-located to reduce communication cost, while
replica anti-affinity and minimum failure-domain rules keep latency
optimisation from defeating availability. Node capacity, architecture,
devices, policy domains, and actual load all constrain the decision.

Routing follows the same abstraction. Clients address a logical service
through the cluster name service. The runtime may choose local delivery when caller and
callee are co-located and network delivery when they are not. If placement
changes, the name's generation changes, stale endpoints are invalidated, and
clients resolve again.

Migration is consequently a cluster operation. Stateless reassignment starts an equivalent instance at
the new location and retires the old one under generation fencing. Stateful
migration additionally requires an explicit transfer of owned mutable state
and a single, fenced change of ownership. CharlotteOS has demonstrated only
the stateless cross-node case; the stateful protocol remains a target.

\section{Kubernetes and OpenShift as a counterpoint}

\subsection{What the container stack constructs}

A Linux container is not a new kernel object equivalent to a virtual
machine. It is a group of ordinary Linux processes made into a deployable
unit by combining several facilities: namespaces determine which resources
and identifiers are visible, cgroups account for and limit resource use,
Linux capabilities divide some root privileges, seccomp restricts the system
call interface, an LSM such as SELinux constrains object access, and a prepared
filesystem and network configuration provide the process's apparent
environment. An OCI image and runtime turn those facilities into a
repeatable execution contract. Kubernetes adds scheduling, desired-state
reconciliation, service naming, failure recovery, rolling change, and many
other cluster functions. OpenShift adds an integrated security, policy,
build, registry, identity, networking, and operations environment around
that model.

This machinery exists for good reasons. Linux is a general-purpose,
machine-centred operating system with broad compatibility obligations.
Conventional applications expect filesystems, shared libraries, package
contents, users, process APIs, network interfaces, and often a distribution
userspace. Container platforms must package enough of that environment to
make a program reproducible, then construct boundaries that let mutually
distrustful workloads share a kernel. OCI images need not contain a complete
distribution, but they commonly carry a substantial userspace precisely
because existing software was written as if it ran inside a machine.

Kubernetes can therefore be understood partly as a distributed
operating-system layer over machines whose native abstractions stop at the
node. It translates the operator's model of cluster workloads into Linux
processes on machines, and gains enormous practical value from orchestrating
the software and operating systems that already exist.

CharlotteOS asks a different question: what becomes possible when the
controlled distributed execution environment is the operating system,
rather than something assembled above a set of general-purpose hosts?

\subsection{Two sources of complexity}

The comparison is useful only if it separates two kinds of complexity.

\subsubsection{Distributed-systems complexity is intrinsic}
CharlotteOS retains the intrinsic responsibilities of distributed systems:
scheduling, consensus, placement, replication, health management, service
discovery, upgrades, resource accounting, storage, networking, observability,
failure detection, load balancing, and recovery from partial failure. Network
partitions require consistency rules; persistent state requires ownership,
durability, repair, and backup; rolling change requires compatibility and
rollback; and secrets require secure creation and delegation. Moving these
responsibilities into the operating system changes their integration, not
their necessity.

\subsubsection{Compatibility and boundary-construction complexity may change}
Some machinery exists because a cluster platform must turn a
general-purpose Unix/Linux environment into a constrained, reproducible,
movable unit. A clean-break system with signed, self-contained artifacts,
separate address spaces, explicit capabilities, cluster-native names, and no
per-node application installation may eliminate or radically simplify:

\begin{itemize}
    \item OCI filesystem images as the normal application-delivery unit;
    \item CRI/OCI container-runtime and low-level container-launch layers;
    \item per-workload namespace construction and much container-specific
        seccomp, Linux-capability, and LSM policy scaffolding;
    \item the attempt to remove ambient machine authority from a process that
        began in a conventional Unix environment;
    \item per-node application package management, configuration drift, and
        machine-centred deployment procedures; and
    \item host-oriented workload mobility assembled from process, filesystem,
        and network emulation.
\end{itemize}

These are hypotheses about mechanisms and their location. Resource accounting,
memory protection, and policy enforcement remain essential even when Charlotte
expresses them through native resource objects, address spaces, and capability
possession instead of cgroups, Linux namespaces, and SELinux labels.

\subsection{The control plane may survive the execution stack}

CharlotteOS explores absorbing cluster orchestration into the operating
system and shortening the machinery required to manufacture a controlled
execution environment above Linux.

Kubernetes's strongest control-plane ideas may survive remarkably well:
declarative desired state, reconciliation, scheduling from constraints,
generation-aware updates, health-driven replacement, replicated control
state, and separation between specification and observation. A CharlotteOS
cluster executive may therefore look recognisably Kubernetes-like at the
level of intent and convergence.

The execution path beneath that controller could, however, be much shorter.
Instead of a desired workload becoming a Pod, an OCI image, a container
runtime request, a low-level runtime invocation, a namespace/cgroup/LSM
configuration, and finally Linux processes, a cluster assignment can become
a verified Charlotte execution object loaded directly into a
capability-constrained address space. Discovering that reconciliation and
scheduling survive while Pods, OCI, CRI, and much container security
scaffolding do not would identify which platform ideas are consequences of
distributed computation and which are consequences of the substrate.

The price of the clean break is equally important: existing Linux software
does not simply run. CharlotteOS exchanges compatibility for a smaller and
more explicit execution contract. Kubernetes asks how to orchestrate today's
applications. CharlotteOS asks what applications and operating systems could
look like if cluster execution, artifact trust, and explicit authority were
native from the beginning.

\section{Target architecture}

\subsection{The minimal node}

The production target is an appliance-like node with no shell, no per-node
application package manager, and no manually curated repository of installed
software. Its stable local environment is just enough to boot, establish
trust and identity, reach storage and the network, join cluster state,
observe resources, and load assigned artifacts. Application differences
belong in signed workload objects and replicated cluster policy rather than
in mutable node configuration.

A target node lifecycle is:

\begin{enumerate}
    \item boot a known CharlotteOS image and establish the node's execution
        and hardware trust boundary;
    \item discover a cluster or contact a configured seed;
    \item authenticate, learn cluster identity and policy, and join the
        relevant replicated state;
    \item advertise resources and failure-domain properties;
    \item receive assignments, fetch immutable artifacts, validate them, and
        instantiate only their delegated authority; and
    \item on failure or retirement, surrender ownership under fencing while
        the cluster reconciles desired state elsewhere.
\end{enumerate}

Secure Boot, measured boot, or remote attestation can extend the
hardware-to-cluster trust chain. Artifact signing establishes executable
identity; node attestation remains future work.

\subsection{The cluster executive}

The target cluster executive holds desired software state and observes
actual execution state. It knows accepted artifact identities, placement
policies, node resources, assignments, active service generations, and
enough health and communication data to converge the two. Its responsibilities
include:

\begin{itemize}
    \item admitting nodes and maintaining consensus membership;
    \item accepting authorised, generation-fenced desired-state changes;
    \item ensuring that referenced artifacts are available and valid;
    \item choosing placements from resource, trust, affinity,
        anti-affinity, and failure-domain constraints;
    \item realising replica counts and routing among active owners;
    \item observing failure and load and initiating replacement or movement;
    \item coordinating upgrades, rollback, and mutable-state ownership; and
    \item exposing cluster-wide status, accounting, logs, metrics, and
        administrative audit.
\end{itemize}

This is substantial orchestration work. CharlotteOS seeks a shorter control
path that acts directly on native execution objects, removing the intermediate
machine-shaped compatibility environment while retaining the hard distributed
systems problems.

\subsection{Placement, routing, and migration as native operations}

The first placement controller can begin with declared rules. Desired
replicas, maximum instances per node, minimum distinct nodes, every eligible
node, affinity groups, anti-affinity groups, required resources, and policy
domains form a constraint problem over advertised node capacity. Admission
must reject combinations the artifact did not authorise: in particular,
concurrent replicas are invalid unless signed metadata blesses parallel
instances.

Declared affinity is only an initial estimate. The runtime should aggregate
actual message counts, byte flow, queue pressure, and executor load across
the cluster. Sustained evidence can then influence placement, just as the
per-node scheduler's sustained-window rebalancing replaces a static
assignment with a load-observed one. Communication-aware placement remains
subordinate to correctness constraints: co-location must not collapse
required failure diversity, violate trust boundaries, or create ambiguous
state ownership.

Routing is the indirection that makes placement change tolerable. A logical
name resolves to an active generation and one or more owners. A local owner
may be reached through a local endpoint and a remote owner through cluster
transport without changing the caller's authority. When an owner is
replaced, generation fencing prevents the retiring instance from deleting or
overwriting its successor's registration.

Stateful migration is explicit. The old instance exports only the
mutable state it owns; the target accepts that state under a transfer
protocol; ownership changes exactly once; the new endpoint becomes active;
and stale clients are forced to re-resolve. Every transfer therefore carries a
consistency and fencing model.

\section{Implemented foundation}

The following implemented and verified mechanisms make the cluster vision
concrete. The remaining controller responsibilities are collected later in
this chapter.

\begin{itemize}
    \item \textbf{Cross-machine consensus.} The Graft core implements generic
        voter sets, learners, quorum decisions, and joint-consensus membership.
        Unit tests exercise configurations including five voters; TLA+ models
        explore three-node safety; and operational QEMU fixtures replicate the
        \texttt{dns} catalog over Ethernet with up to three guests
        (\autoref{chap:consensus}). The catalog carries membership, names,
        deployments, the set-once cluster signing key, signed placement rules,
        concrete replica sets, and per-node readiness.

    \item \textbf{Reliable cluster messaging.} Protocol v3 represents message
        lengths and fragment offsets with 32 bits and initially admits messages
        up to a bounded 1~MiB policy ceiling. They are fragmented across
        Ethernet frames and reassembled at the receiver
        (\autoref{chap:networking}). Remote name-service calls, deployment coordination,
        and the demonstrated assignment handoff use this transport.

    \item \textbf{Persistent per-node object stores.} The NVMe-backed object
        store uses generation-selected directory records and copy-on-write
        replacement (\autoref{chap:storage}) for service artifacts and persistent Raft
        state. The current image producer places the same signed bundle on
        each node for bootstrap. Signed application deployments may instead
        name an opaque key in a centrally managed S3 store; the assigned node
        pulls through its locally provisioned connector. Repair and
        cluster-wide artifact replication are not implemented.

    \item \textbf{Capability-scoped external data planes.} Separately
        provisioned S3 and Kafka services run above the userspace TCP/IP and
        verified-TLS path (\autoref{sec:external-data-services}). An S3
        capability confines an application to one managed endpoint, bucket,
        prefix, credential identity, and rights mask. A Kafka capability
        confines it to an explicit broker-destination allow-list, fixed consume partition, bounded
        produce-route allow-list, group, transactional identity, and operation
        mask while keeping TLS/SCRAM and mTLS credentials inside the connector and offering idempotent production,
        read-committed consumption, and transactional offset commits. These
        adapters make externally managed infrastructure usable through
        Charlotte's capability model. Cluster-native storage and messaging are
        separate research questions.

    \item \textbf{Role-separated operational deployment.} Development signs
        an immutable release and its deployment descriptors, while operations
        independently signs HPKE-encrypted S3 or Kafka connector profiles bound
        to that exact release, target connector, cluster, expiry, and central
        object-store key. The leader verifies both authorities and trusted UTC,
        then places only compact references and replay fences in Raft. On the
        assigned node, the authorized agent retrieves the digest-pinned
        ciphertext through a bootstrap S3 capability and moves a bounded pickup
        object to the kernel. The kernel repeats release, descriptor, artifact,
        envelope, and expiry verification, decrypts into transient zeroizing
        memory, validates the connector-specific profile, and transfers it as a
        read-only launch capability. Application IPC receives only attenuated
        connector endpoints, never infrastructure credentials. Production key
        custody, rotation, audit, and readiness-driven generation cutover remain
        future work.

    \item \textbf{Local stateful replacement.} Live service upgrade
        transfers a service's state as a memory object, tracks generations,
        invalidates stale connections, and re-registers the replacement
        through the name service (\autoref{chap:live-upgrade}). This has not yet been
        extended to mutable state crossing the network.

    \item \textbf{Service isolation and routing.} Services execute in
        separate address spaces with delegated capabilities
        (\autoref{chap:service-architecture}).
        The name service is replicated across nodes and provides the routing
        indirection for a service somewhere in the cluster.

    \item \textbf{Certified work migration.} The scheduler migrates threads
        between LPs with generation-validated handoff and sustained-window
        load rebalancing (\autoref{chap:scheduling}). This per-node mechanism
        supplies design evidence for future cluster-wide load rebalancing.

    \item \textbf{Observation substrate.} The sitas executor exposes task
        counts, poll counts, queue depth, and shard snapshots
        (\autoref{chap:sitas-executor}).
        Shard-to-shard message flow can be measured. Cluster-wide collection
        and use of those measurements by a placement controller remain
        future work.

    \item \textbf{Signed binaries at the loading boundary.} The EL0 loader
        validates ELF structure, rejects writable-executable segments, maps
        images with page-level permissions, and requires an Ed25519 CLS2 note
        bound to the requested service name (\autoref{chap:storage} and
        \autoref{chap:live-upgrade}). Store
        lookup, administration ingress, and deployment recheck that identity.
\end{itemize}

\architecturewidefigure{external-service-capabilities}
    {Secret-bearing S3 and Kafka connectors receive infrastructure profiles;
    applications receive only named, role-attenuated service capabilities.}
    {fig:external-service-capabilities}

\architecturefigure{operational-profile-pickup}
    {Privileged operational-profile pickup keeps ciphertext references in Raft
    and decrypts directly into immutable connector launch memory.}
    {fig:operational-profile-pickup}

\section{Current bootstrap and node model}

The implemented kernel embeds five signed bootstrap services:
\texttt{ns}, \texttt{nvme}, \texttt{objstore}, \texttt{uart}, and
\texttt{observe}. They establish naming, reach persistent storage, and
provide the early console and observation paths. Other staged services are
loaded on first use from the signed bundle in the initial NVMe object-store
image and cached in kernel memory. Service binaries therefore no longer
generally need to be linked into the kernel image.

This establishes the minimal-node bootstrap shape. The current
\texttt{--deploy-test} serial command loop remains a kernel-side validation and
administration shim; production management enters through userspace deployment
ingress and still needs further hardening.

\subsection{Discovery and membership}

The operational DNS-owned Raft group now admits members discovered on the
cluster Ethernet through replicated \texttt{JOIN}, joint consensus, and
\texttt{FINALIZE}. The three-guest ingress fixture exercises that path before
driving placement and VIP failover. A smaller two-guest fixture remains useful
for focused catalog and deployment checks; its topology is not a membership
restriction.

A production node joining shared cluster state must complete three distinct
operations:

\begin{enumerate}
    \item \textbf{Discover the cluster.} The node uses broadcast discovery
        or a configured seed as discussed for Raft peers in
        \autoref{chap:consensus}.
        Discovery exists as a service, supplies MAC routes and cluster
        posture, and locates the anchor in the generic admission test.
        Discovery alone does not grant voting authority.

    \item \textbf{Receive cluster identity.} The replicated catalog can
        distribute the set-once public signing key. Today that key must
        equal the build-time bootstrap anchor. Establishing trust from a
        genuinely blank start and rotating keys are not implemented.

    \item \textbf{Join shared state.} The node becomes a member of the
        replicated state machine and begins receiving assignments. The leader
        resolves signed singleton, fixed-replica, every-eligible-node, affinity,
        and anti-affinity policy into concrete node sets and reconciles them
        after membership or drain changes.
\end{enumerate}

The target may allow a node to discover peers and participate in a restricted
bootstrap before cluster key material is established, but it cannot validate
executable artifacts without an accepted public key. An authorised
administrative operation would establish or rotate that trust through
replicated state. The current prototype instead uses a build-time anchor and
a set-once, idempotent key ceremony.

\section{Software ingress and deployment}

\subsection{The target pipeline}

CharlotteOS gives the software supply chain an explicit cluster boundary. The
intended pipeline is:

\begin{enumerate}
    \item build a self-contained workload from reviewed source and
        dependencies;
    \item evaluate policy and retain provenance, SBOM, test, and attestation
        evidence;
    \item bless and sign the immutable workload identity off-cluster;
    \item make the signed artifact available to the cluster; and
    \item change replicated desired state so the cluster validates, places,
        and runs the accepted version.
\end{enumerate}

The signing private key remains off-cluster, and placement fetches complete
artifacts instead of executable dependencies from the Internet. The same
immutable object can consequently run on any eligible node without
reconstructing a userspace installation there.

\architecturefigure{release-admission-rollout}
    {A signed multi-component release is admitted atomically to desired state;
    artifact retrieval, launch, publication, and readiness then reconcile per
    component.}
    {fig:release-admission-rollout}

\subsection{The implemented pipeline}

The current deployment pipeline has three stages:

\begin{enumerate}
    \item \textbf{Blessing and signing.} Version-controlled policy admits an
        artifact and the host-side build and \texttt{cluster-sign} tooling
        sign it with the cluster key. CLS2 metadata binds the logical name,
        class, release, rollback counter, parallel-instance permission, and
        an optional SHA-256 link to retained SBOM/build-provenance evidence.
        Nodes receive only the public key.

    \item \textbf{Validation.} Each node validates every loaded ELF against
        the public bootstrap anchor and, in the deployment path, the matching
        replicated key. The loader, object-store resolver, administration
        ingress, and agent enforce signature, logical-name, and digest checks.
        The present ceremony can distribute only the build-time anchor; it is
        not a general enrolment or rotation protocol.

    \item \textbf{Assignment.} A Raft-committed record directs a concrete
        replica set to load a particular object id and immutable ELF digest.
        Each selected node's agent fetches the opaque object key through its
        locally provisioned S3 capability, then validates and runs those exact
        bytes. Endpoint and credential details remain inside that privileged
        connector rather than the descriptor or application.
\end{enumerate}

The replicated catalog records a sorted concrete node set per artifact and an
independently generation-fenced readiness registration for each active owner.
The leader resolves \emph{somewhere}, fixed replicas, and
\emph{every eligible node} against admitted, non-draining voters and
reconciles the assignments when that set changes. DSR routes new cluster-VIP
flows across the ready subset. The boundary between this management loop and
the DSR packet path is shown in
Figure~\ref{fig:cluster-management-dsr}: placement and readiness determine the
eligible set, while each packet uses an already materialized snapshot.

\section{Placement: contract, observation, and movement}

\subsection{Declared policy and bounded scheduler}

The shared \texttt{PlacementPolicy} type expresses desired replicas, maximum
instances per node, minimum distinct nodes, \emph{every eligible node}, an
affinity group, and an anti-affinity group. A request for concurrent
replicas is invalid unless the artifact's signed CLS2 metadata contains
\texttt{PARALLEL\_INSTANCES}. \texttt{CDEPLOY5} signs this policy, the
leader resolves it to concrete assignments, and catalog-v13 snapshots retain
both the desired set and per-node readiness. Affinity groups share a stable
ranking seed; anti-affinity groups require unused nodes and fail closed when
the eligible set cannot satisfy the separation. A policy that needs more
distinct nodes than the current eligible set fails admission.

The scheduler remains intentionally bounded. It does not yet account for CPU
or memory capacity, labels, named failure domains, or more than one instance
of an artifact on a node. Reassignment advances the whole deployment
generation, so retained replicas restart rather than participating in a
rolling surge/unavailability protocol.

\subsection{Runtime observation: local substrate only}

Declared relationships are guesses. Shard-level message counts, byte flow,
queue depth, and task snapshots are already observable inside a node
(\autoref{chap:sitas-executor}), and the distributed name service can identify which components
call which. The target controller should aggregate these signals across
nodes and use sustained behaviour, not transient noise, to refine placement.

No cluster-wide metrics aggregation or feedback-driven placement loop is
implemented today. The existing executor measurements are inputs from which
that controller can later be built.

\subsection{Focused stateless-handoff validation}

The focused two-guest fixture starts the signed \texttt{greet} artifact on one
node, invokes it remotely, reassigns it, retires the old instance, and verifies
that the logical name remains reachable. Generation fencing prevents the
retiring owner from unregistering its replacement. This fixture isolates
assignment handoff and routing continuity for fresh service state; load-driven
and stateful migration require additional protocols described below.

\subsection{Future stateful cross-node migration}

Stateful migration must extend the local live-upgrade protocol from
\autoref{chap:live-upgrade}:

\begin{enumerate}
    \item the cluster instructs the target node to obtain, validate, and
        prepare the component in a new address space;
    \item the old instance exports explicitly owned mutable state and the
        cluster transfers that state to the prepared instance;
    \item the target accepts ownership under a fencing protocol;
    \item the new instance registers through the distributed name service,
        advancing the name generation;
    \item clients resolving the name reach the new owner, while clients with
        stale connections observe \texttt{EndpointClosed} and re-resolve; and
    \item after the new owner is committed and reachable, the old instance is
        retired and its resources reclaimed.
\end{enumerate}

Generation tracking, endpoint invalidation, name-service re-registration,
and local memory-object state transfer are verified primitives. Network
transport and ownership fencing for mutable service state are not
implemented. The demonstrated cluster handoff starts the same immutable
artifact with fresh state.

\section{Deployment, rollout, and reassignment}

The deployment-test QEMU flow uses two guests to isolate one source-to-target
handoff. A signed artifact is assigned through replicated state; the remote
agent obtains and verifies it; the service is invoked across the network; and
the assignment moves without losing the logical name. Both guests complete
with zero failed or pending self-tests. The operational membership, replica,
and ingress paths also run in a three-guest fixture, and the underlying Raft
membership is generic. These fixture sizes describe test coverage only.

\architecturewidefigure{two-node-cluster}
    {A focused two-guest deployment fixture combines Raft-held desired state,
    local reconcilers, capability-routed services, and central object pickup;
    the architecture admits larger voter and replica sets.}
    {fig:two-node-cluster}

\subsection{Replicated deployment manifest}

The distributed name service's Raft catalog (\autoref{chap:consensus}) also holds a
\textbf{deployment manifest}:

\begin{quote}
\texttt{artifact -> \{object identity, artifact digest, replica nodes,
generation, signed descriptor\}}.
\end{quote}

The map is committed through the same log as name registrations. Deployment
is therefore ordered, replicated, generation-fenced cluster state rather than
per-node configuration. \texttt{OP\_DEPLOY} submits a record and
\texttt{OP\_DEPLOY\_QUERY} reads a replica's applied state.

The signed descriptor authorises the deployment decision and binds artifact
identity, placement policy, and resource limits. Raft orders and replicates the
result. Production ingress still needs a separately delegated online
administrator identity and audit policy in addition to the offline signature.

\subsection{Cross-node hosting and generation fencing}

A service hosted away from the leader registers through a
\textbf{remote-register relay}, because only the leader may commit catalog
entries. The leader returns the committed generation. The two-phase
register/activate sequence and generation-fenced endpoint-death unregister
are preserved across the relay, so a retiring host cannot clobber a
replacement's registration.

Cross-node invocation uses the existing \texttt{OP\_CALL} relay. A client on
any member reaches the service at its current owner through the replicated
name service rather than through a node-specific application address.

\subsection{The deployment agent}

Each networked node with durable local storage runs an \texttt{agent}. It polls
the manifest for assignments to its node key. A legacy record selects bytes
from the local object store. A signed descriptor instead names an opaque key
that the agent pulls through the separately provisioned S3 connector. In both
cases it verifies the assigned SHA-256 identity and checks that the signature
note covers both the bytes and requested logical name. The descriptor path
then invokes the scoped deployment-authority syscall with both owned memory
objects.

The kernel snapshots and revalidates the exact ELF and descriptor, maps the
ELF into a fresh address space, gives the application only the grant controller
and its immutable descriptor, and starts it. A descriptor-authorized service
publishes its endpoint through \texttt{grantctl}; the agent supervises the
assignment. When the manifest assigns the artifact elsewhere, the agent
retires the service and causes its address space to exit. Endpoint closure is
observed by \texttt{dns}, whose generation check prevents a stale unregister.

The supervisor-designated agent authority is the one narrow kernel mechanism:
only that agent may start or retire a verified, name-bound ELF memory
object. The desired placement and administration logic remains a userspace
coordination concern.

\subsection{\texttt{clusterctl} and the test console}

\texttt{clusterctl} is the prototype outside administration interface. It
wraps the DNS manifest operations and object store behind:

\begin{itemize}
    \item \texttt{upload}, which stores a note-signed artifact under its
        derived cluster-wide id after checking that CLS2 blesses the requested
        name;
    \item \texttt{deploy}, which hashes the stored bytes and submits the
        pinned assignment; and
    \item \texttt{status}, which queries the committed record; and
    \item \texttt{notify}, which verifies a signed \texttt{CDEPLOY5}
        descriptor and commits its opaque central-object key, digest, target,
        sequence, per-thread stack pages, maximum active-thread count, and
        cooperative shutdown grace period and capability grants without
        copying an ELF through Raft.
\end{itemize}

The normal service set also starts \texttt{deployd}. Its bounded plaintext
HTTP endpoint accepts \texttt{POST /v1/deployments} on guest port 7444 with
exactly one signed descriptor as the body. The signature is the authorisation
and integrity envelope and the descriptor carries no secret. Lower sequences
are rejected, different bytes at an existing sequence conflict, and identical
notifications are idempotent. QEMU's user network forwards host port 8081 by
default; \texttt{cluster-sign deployment-notify} is the corresponding host
client.

For a release made from several independently signed components,
\texttt{cluster-sign deployment-apply} prevalidates all descriptor files,
rejects duplicate artifact names, submits each descriptor, and waits for every
exact generation to become ready against one deadline. Artifacts must already
be signed and uploaded; object-store credentials remain outside the generated
application and its descriptors.

\texttt{deployment-apply} orchestrates independent descriptor admissions.
Each descriptor occupies a separate Raft entry, so a later failure can leave
an accepted prefix committed. Identical retries are idempotent. The atomic
alternative is described next.

A signed \texttt{CRELEASE} envelope produced with
\texttt{cluster-sign release-sign} and submitted with
\texttt{release-apply}. Its outer signature binds a release name, monotonic
sequence, and the exact ordered bytes of up to 16 distinct nested descriptors;
the complete envelope is bounded to 3,584 bytes. The cluster verifies both the
outer signature and every nested deployment signature. Current tooling emits
\texttt{CDEPLOY5}; the reader retains compatibility with \texttt{CDEPLOY1}
through \texttt{CDEPLOY4}.
A follower relays the
request to the current leader, which resolves all signed placement policies
against admitted, non-draining voters and proposes the concrete replica sets
as one Raft command. The state machine preflights every
release and component revision before changing the deployment map, so desired
state changes all at once or not at all. Exact retries are idempotent and the
signed release record survives catalog snapshots.

After atomic admission, agents fetch, launch, and publish independently;
\texttt{release-apply} waits for every exact generation. Coordinated rollback,
progress policy, and a semantic process bundle that also binds BPMN, schemas,
and provenance remain controller work.

Retirement is nevertheless bounded locally. On the first reconciliation poll,
the kernel publishes a drain request and monotonic deadline in the domain's
read-only launch page. A lifecycle-aware application stops new work, completes
or aborts external transactions, returns from its resource-owning scope, and
then exits. If any domain thread remains, the kernel forces termination on the
agent's first retirement poll at or after the signed deadline and reclaims the
address space only after generation-safe reaping. While retirement is active,
the agent polls at most ten milliseconds apart.
Figure~\ref{fig:cooperative-deployment-shutdown} shows this path.

Production whole-node shutdown begins off-cluster as a fixed-size, signed
\texttt{CSHUTDN1} intent. It binds one nonzero target node, a per-node
monotonic sequence, a UTC validity interval, whole-node and per-phase grace
budgets, and the power-off reason. There is deliberately no broadcast target:
an operations controller can drain nodes in an order that preserves Raft and
application quorum. The operator creates the record with
\texttt{cluster-sign shutdown-sign} and submits it to
\texttt{POST /v1/shutdowns}; no private key enters CharlotteOS.

Any member may receive the bounded envelope. A follower relays it over the
reliable-message channel, while the leader verifies the committed cluster key
and trusted UTC before proposing it. The catalog state machine verifies the
signature again, fences stale and conflicting per-node sequences, and retains
the exact envelope in snapshots. The target agent consumes only locally
applied state, re-verifies the envelope, and moves its memory capability into
the kernel. The kernel admits only its designated deployment agent,
independently checks the signature and target node, derives a local monotonic
deadline from the signed duration, and transfers the service graph to a kernel
worker. Thus retiring the agent cannot orphan a partially completed node
drain.

The same lifecycle mechanism then propagates the node drain through the deployment plane.
On receiving a kernel-authenticated \texttt{NodeShutdown} lifecycle request,
the agent stops reconciliation and launches no new work. It marks every
ordinary deployment and privileged operational connector as retiring, then
retains and polls their owners until reclamation completes. Each child receives
the earlier of its descriptor-signed grace deadline and the enclosing node
deadline: the node transition may shorten an allowance but cannot extend it.
Ordinary applications drain before their privileged operational connectors,
so an application retains Kafka/S3 authority while it finishes or aborts
external work. The agent acknowledges and exits only after both child sets are
empty.

Deployment ingress, HTTP ingress, and the UTC time service also implement the
cooperative path. The ingress services stop opening listeners and poll their
current accept or receive operation with a bounded lifecycle check before
dropping their owned listener and service connections. A receive timeout uses
an explicit cancellation exchange with TCP/IP: closing only the local call
capability would leave the server's deferred reply slot occupied and could
discard the first bytes that arrive later. Time closes its public endpoint and
unwinds any pending NTP call, UDP socket, and persistence connection as one
owning scope. Already admitted ingress work may finish inside the signed grace,
but none of these services begins new work after observing the lifecycle
request.

The transport phase is cooperative too. TCP/IP runs after every socket
consumer has drained; it closes admission, completes retained receive calls,
stops residual protocol sockets, and releases its endpoint and owned NIC
connection. The frame router then drops one owning scope containing the single
NIC receive, deferred name lookups, routed service connections, and in-flight
moved-frame calls. The NIC driver cannot transfer to device quiescence until
both transport domains have exited.

Startup gates obey the same lifecycle contract. Long-lived services poll the
node-ready name through bounded non-blocking lookups, checking the authenticated
lifecycle page between attempts. HTTP, time, and TCP/IP can consequently
acknowledge a node drain even if shutdown begins before local readiness is
published. Before reclaiming a domain, the coordinator records per-phase
counts for acknowledged, unacknowledged, and forced exits.

The deployment plane and cluster core participate in cooperative shutdown.
Deployment ingress closes its listener, cluster
control drops its endpoint and service connections, and the agent interrupts
startup or reconciliation waits before draining applications and then their
operational connectors. DNS stops the Raft/catalog reactor and cancels pending
transport and catalog work. Reliable messaging rejects retained sends,
completes its deferred receive, and releases queued messages before discovery
stops probing and settles pending cluster-status waiters. Bounded
lifecycle-aware sleeps ensure a configured retry interval does not become
shutdown latency.

The kernel's next shutdown layer transfers the published steady-state handles
out of the launch registry and retires high-level node services in reverse
dependency order: deployment ingress, deployment control, deployment agent;
HTTP and time; cluster catalog, reliable messaging, and discovery; TCP/IP; and
the frame router; and finally the object store. Each phase must be
generation-safely reclaimed before the next receives its request. Per-phase
grace is capped by the enclosing node deadline.

The remaining NIC, block, and entropy domains transfer exactly once to a
device-quiescence coordinator. It never converts a deadline into an
unconditional driver abort. Ordinary lifecycle readiness is insufficient: a
driver must publish the distinct \texttt{DEVICE\_QUIESCED} state and exit before
the kernel reclaims its address space and launch-owned device grants. Otherwise
the domain and IOMMU state remain quarantined for diagnosis. Every in-tree
hardware adapter implements the driver side of this contract. VirtIO entropy
resets the device before dropping its shared-DMA mappings. NVMe
adapter closes admission, drains outstanding replies and DMA, observes a final
NVM Flush completion, masks interrupts, clears \texttt{CC.EN}, and waits for
\texttt{CSTS.RDY=0}. VirtIO block flushes before device reset; AHCI flushes,
masks port interrupts, and stops its command and FIS receive engines; VirtIO
net and E1000E close ingress and deferred receives, release queued frames,
drain transmit descriptors, and then reset or halt the receive/transmit
engines. A failed flush or drain still stops DMA where possible but withholds
\texttt{DEVICE\_QUIESCED}. Thread death alone is never treated as evidence
that DMA stopped.

After verified device quiescence, AArch64 selects the FADT-advertised PSCI SMC
or HVC conduit, masks interrupts, and invokes \texttt{SYSTEM\_OFF}. It must not
return: scheduling cannot resume past terminal drain. The fixture accepts QEMU
exit only after the authoritative success and poweroff records.

\architecturewidefigure{cooperative-deployment-shutdown}
    {Bounded cooperative shutdown for deployment retirement or an enclosing
    node drain, followed by generation-safe forced retirement. The application
    can acknowledge but cannot clear or extend a kernel-published lifecycle
    request.}
    {fig:cooperative-deployment-shutdown}

The client may contact any member. A follower relays the bounded submission to
the current Raft leader over a source-validated peer route and correlates the
committed result back to the ingress request. The management client therefore
does not need a separate leader-discovery protocol.

The same endpoint accepts
\texttt{GET /v1/deployments/\{percent-encoded-name\}}. It reports
\texttt{committed}, \texttt{replacing}, or \texttt{ready}, together with
desired and ready replica counts. A ready observation is generation-safe: all
desired nodes must have active registrations carrying the exact deployment
generation. A stale endpoint with the same name cannot make a replacement
rollout appear ready.

The signed descriptor may carry node key zero to request automatic placement.
A singleton prefers the current Raft leader and uses a deterministic fallback
when that node is draining; fixed replica counts and
\emph{every eligible node} select distinct admitted, non-draining voters. The
leader continuously reconciles that set after membership changes. A nonzero
key remains an explicit singleton pin. Replica policy is admitted through a
\texttt{CRELEASE}, giving affinity decisions the complete component context.
Full signed artifact names
up to the 48-byte CLS2 limit travel in memory-backed catalog operations rather
than the historical eight-byte scalar encoding.

An interactive kernel serial console in the \texttt{--deploy-test} flow issues
these commands against the live cluster for validation. Production operations
use the userspace notification endpoint, whose remaining hardening includes
network-level denial-of-service controls, an online audit identity distinct
from the signing key, and production transport policy. The raw mutation path
and serial shim remain test-only privileged interfaces.

The node agent now reconciles the complete replicated deployment-name set,
rather than polling only \texttt{greet}. A node can own up to 64 independently
named application domains. Each domain is fenced by the principal derived from
its signed artifact name, so a generation change retires only the matching
application even if address-space identifiers have been reused. This bound is
kernel admission policy, not a naming or descriptor-format limit.

Automatic placement now covers singletons, fixed distinct-node replica sets,
every eligible voter, stable affinity/anti-affinity ranking, rescheduling, and
exact-generation readiness. The controller commits desired assignments
through Raft while node agents remain mechanism-only reconcilers. It still
needs capacity-aware selection, labels and named failure domains,
multi-instance-per-node execution, rollout surge/unavailability, and rollback.

The development runner's
\texttt{--deployment-ingress-test --timeout 240} option proves the complete
host upload to TLS RustFS, signed atomic release admission, S3 pull, scoped
launch, and generation-safe readiness path. It is an integration fixture;
production still needs a release controller that manages policy approval,
progress deadlines, rollback, and operator audit identity.

\subsection{Artifact naming}

Artifact names are currently bare cluster-global names such as
\texttt{greet}. The compatibility object id is derived from the name using a
\texttt{0xfffe...} tag over FNV-1a, while signed descriptors carry the opaque
central-S3 object key and full content digest. Concrete node assignments occur
only in the replicated deployment record.

The host image producer detects and rejects collisions in the complete
bundle. Production should replace the truncated 48-bit name hash with a
collision-resolving or fully content-addressed namespace.

\subsection{Cryptographic boundary}

\architecturewidefigure{role-separated-deployment-trust}
    {Independent artifact, deployment, operations, and recipient-key roles
    converge only at cluster admission and the privileged connector launch
    gate.}
    {fig:role-separated-deployment-trust}

Artifact validation uses Ed25519 signatures embedded in the ELF. The private
cluster key remains off-cluster in the
\repodir|tools/cluster-sign| host tool. The signature is stored in a standard
ELF \texttt{SHT\_NOTE} section named
\path{.note.charlotte-sig}, with note type \texttt{"CLS2"}. Its descriptor
contains the logical name, artifact class and flags, release, rollback
counter, signing-key id, optional provenance digest, and the 64-byte
signature.

The signature covers the SHA-256 digest of the complete file with only the
signature field treated as zero, making the signed representation
deterministic. The matching public key enters the prototype in two ways: it
is embedded in the kernel at build time as the bootstrap trust anchor and
passed to services in the launch manifest; and the key ceremony commits it
to replicated state through
\texttt{OP\_KEYCEREMONY} and \texttt{OP\_SET\_KEY}. This demonstrates a
distribution path for dynamically joining nodes. The replicated value must
currently equal the build-time anchor; blank-start trust establishment remains
future work.

The EL0 loader verifies the note before mapping a domain. Signing is
mandatory: unsigned or invalidly signed images are rejected. The
live-upgrade syscall reports failure gracefully; other current load paths
panic. By default the build signs staged services using the
version-controlled development keypair at
\repofile|tools/cluster-sign/dev-key.hex|, whose public half is the embedded
\texttt{CLUSTER\_PUBLIC\_KEY}. A different key can be supplied through
\texttt{CLUSTER\_SIGN\_PRIVATE\_KEY}, with a kernel rebuilt to embed its
public half. A production key would be held outside the repository under
organisational key-management procedures.

Store resolution and deployment additionally check the signed logical name,
and the manifest pins the full ELF SHA-256. The demonstrated
\texttt{greet} service is therefore a real, signed binary fetched from the
local object store, mapped, executed, served across the network, and
reassigned. Raft orders the assignment, but consensus is not an
authorisation proof for the administrator who requested it.

\subsection{Validated properties}

The focused deployment fixture validates that:

\begin{itemize}
    \item an explicit node assignment can be represented as replicated,
        generation-fenced cluster state;
    \item a signed, name-bound ELF can be obtained, independently revalidated,
        and launched by a node agent;
    \item a service can be hosted away from the Raft leader and remain
        reachable through the distributed name service;
    \item generation fencing protects a stateless reassignment from a stale
        owner's unregister; and
    \item deployment and administration have protocol shapes rather than
        requiring direct machine configuration.
\end{itemize}

Separate multi-guest fixtures validate dynamic DNS/Raft admission, replica
placement, and DSR failover. Evidence is still required for resource-aware
placement, rolling replacement, stateful cross-node migration, artifact-store
recovery, production key establishment and rotation, and production
administrative authorisation.

\section{Remaining implementation boundary}

The gap between the demonstrated functionality and the target cluster operating
system contains the following substantial coordination and trust work:

\begin{itemize}
    \item \textbf{Dynamic membership recovery.} Discovery-bridged admission
        performs replicated \texttt{JOIN}, joint consensus, and
        \texttt{FINALIZE} in the durable DNS-owned group. Partition behaviour,
        restart and rejoin within a live quorum, automatic failed-voter
        removal, and longer-running fault injection remain.

    \item \textbf{Trust establishment and key rotation.} The present key
        ceremony is set-once, idempotent, and limited to the build-time
        anchor. Authenticated blank-start establishment, overlapping
        rotation, revocation, and recovery remain.

    \item \textbf{Durable artifact availability.} Signed descriptor deployment
        pulls digest-pinned artifacts through a separately provisioned S3
        connector; the legacy path uses per-node NVMe staging. Production work
        includes integrity-checked repair, retention, garbage collection,
        connector bootstrap recovery, and replication policy.

    \item \textbf{Cluster-wide observation.} Per-executor metrics must be
        collected and interpreted across nodes without turning transient
        telemetry into unstable placement decisions.

    \item \textbf{Placement and rollout depth.} The leader currently resolves
        fixed and every-node replica policy, basic affinity/anti-affinity, and
        drain-aware membership into assignments. The next scheduler adds node
        capacity, trust labels, named failure domains, health and communication
        signals, multi-instance-per-node execution, and rolling
        surge/unavailability and rollback policy.

    \item \textbf{Replica health and failure convergence.} Exact-generation
        readiness gates DSR and replica counts converge after committed
        membership or drain changes. Health-driven replacement, automatic
        failed-member removal, and state-aware recovery remain controller work.

    \item \textbf{Stateful cross-node migration.} The stateless handoff must
        be extended with network state transfer, explicit ownership
        acceptance, fencing, failure recovery, and a consistency model.

    \item \textbf{Production administration authority.} The userspace
        \texttt{clusterctl} protocol exists, but the raw mutation endpoint
        and kernel serial shim do not enforce the intended separately
        delegated cluster-administrator capability. Authentication,
        authorisation, audit, and safe concurrent administration remain.

    \item \textbf{Operational cluster services.} Production health
        management, upgrades, resource accounting, persistent storage
        policy, networking policy, secrets, logging, metrics, tracing,
        backup, and disaster recovery remain necessary even where their
        mechanisms differ from Kubernetes/OpenShift.
\end{itemize}

Most of this belongs in userspace coordination over the existing capability,
messaging, storage, and consensus primitives. Future work must retain the
discipline of the current mechanisms: explicit ownership, bounded
communication, immutable artifact identity, generation fencing, and
replicated decisions instead of per-machine mutation.

\section{Related work}

The vision combines ideas from distributed operating systems, capability
systems, cluster schedulers, immutable software stores, and artifact-trust
systems. The comparisons below identify intellectual neighbours; they are
not claims of equivalent guarantees or implementation maturity.

\subsection{Historical distributed operating systems}

\begin{itemize}
    \item \textbf{Amoeba} (Vrije Universiteit Amsterdam) is a close ancestor
        of the processor-pool model. Work can run on processors in a pool,
        dedicated machines provide file and directory services, and
        network-wide RPC uses capability-based access. Amoeba deliberately
        did not migrate processes; explicit cross-node migration is an
        extension in the CharlotteOS target.

    \item \textbf{Plan 9} (Bell Labs) separates terminals, CPU servers, and
        file servers, and centralises authentication and key handling in
        \texttt{factotum}. Its machine roles and network-transparent names
        challenge the assumption that every node is an independent Unix
        installation.

    \item \textbf{Inferno} continues the Plan~9 line with a portable virtual
        machine and code distributed as data over the network.

    \item \textbf{Jini and JavaSpaces} contribute a discovery pattern in
        which services locate a lookup service, register, and provide clients
        with a proxy. JavaSpaces' tuple-space model resembles work or
        software being placed in shared space for an eligible participant to
        take.

    \item \textbf{MOSIX} explores a single-system-image cluster with resource
        discovery and transparent process migration. Its work on
        communication-aware placement, including Keren and Barak,
        \emph{Opportunity Cost Algorithms for Reduction of I/O and
        Interprocess Communication Overhead in a Computing Cluster}, IEEE
        TPDS 2003, is a precedent for using interoperation cost in placement.
\end{itemize}

\subsection{Cluster control planes and placement}

\begin{itemize}
    \item \textbf{Borg, Omega, and Kubernetes} established the modern
        cluster-management shape: desired software is submitted to a pool of
        machines, scheduling is a cluster decision made from constraints and
        packing, and controllers reconcile observation with desired state.
        Burns, Grant, Oppenheimer, Brewer, and Wilkes describe this lineage in
        \emph{Borg, Omega, and Kubernetes: Lessons from Three Decades of
        Platform-as-a-Service}, ACM Queue 2016.

    \item \textbf{Kubernetes and OpenShift} are also the central counterpoint
        of this chapter. CharlotteOS is not attempting to wish away their
        distributed control responsibilities; it is testing whether a native
        execution and authority model can substantially shorten the machinery
        underneath them.

    \item \textbf{HashiCorp Nomad and Consul} pair cluster scheduling with
        service discovery, membership, and affinity or anti-affinity
        constraints.

    \item \textbf{VMware DRS and vMotion} demonstrate the full operational
        arc of declared placement rules, load-observed rebalancing, live
        movement, and network-level switchover, though for virtual machines
        rather than Charlotte execution objects.

    \item \textbf{Erlang/OTP distribution} provides node discovery, named
        processes, supervision, hot code loading, and patterns resembling
        load, register, redirect, and retire.
\end{itemize}

\subsection{Capabilities and constrained execution}

CharlotteOS also belongs in the lineage of capability-oriented and
least-authority systems. The relevant distinction is between possessing an
explicit reference that confers authority and inheriting a broad machine
environment whose authority must be denied by policy. Capability systems,
microkernels, language sandboxes, WebAssembly runtimes, and unikernels make
different trade-offs along this boundary. CharlotteOS combines a
capability-constrained service model with signed native artifacts and a
cluster-level placement objective; similarity in one dimension does not
imply identical isolation or compatibility guarantees in the others.

\subsection{Trust, signing, and bootstrap}

\begin{itemize}
    \item \textbf{The Update Framework (TUF)} defines signed, versioned
        metadata and key roles for secure update systems. Its separation of
        root trust, online roles, expiration, and rotation is relevant to the
        blank-start and key-lifecycle work not yet implemented.

    \item \textbf{Uptane} applies secure-update principles in the automotive
        setting and illustrates offline and provisioned-key patterns.

    \item \textbf{Sigstore/Cosign and Notary} provide related approaches to
        signing, identity, provenance, and policy for deployable artifacts.

    \item \textbf{Measured boot and remote attestation}, including TPM-,
        ARM CCA-, and confidential-VM-based mechanisms, can complement
        artifact validation by establishing properties of the node on which
        the artifact will run.
\end{itemize}

\subsection{Membership and artifact stores}

\begin{itemize}
    \item \textbf{SWIM} (Gupta, Birman, and van Renesse, DSN 2002) and
        implementations such as \textbf{memberlist} provide scalable
        membership and failure-detection techniques relevant to discovery.
        Failure detection and Raft voting admission remain distinct
        responsibilities.

    \item \textbf{Nix} demonstrates immutable, hash-derived software-store
        paths and reproducible dependency closure, closely related to the
        goal that software be an identified object rather than a mutation of
        a node.

    \item \textbf{OCI registries with signed manifests} demonstrate
        content-addressed distribution and integrity for today's container
        delivery model. A CharlotteOS artifact service needs analogous
        availability and lifecycle properties even if it does not store OCI
        filesystem images.
\end{itemize}

\section{Research thesis}

The cluster vision can be summarised as the following proposition:

\begin{quote}
    Containers are partly an artifact of retrofitting isolation,
    reproducibility, software distribution, and workload mobility onto
    general-purpose operating systems whose primary administrative boundary
    remains the machine. If immutable artifact identity, capability-based
    authority, constrained execution, cluster-wide naming, and placement were
    native operating-system abstractions, much of the contemporary container
    execution stack could become unnecessary, while the essential
    distributed control plane would remain.
\end{quote}

CharlotteOS is an experimental vehicle for testing that proposition. Evidence
now spans signed and name-bound loading, capability-separated services,
replicated naming and deployment state, dynamic membership, deterministic
replica placement, central artifact pickup, remote invocation,
generation-safe stateless handoff, and cluster-wide TCP ingress. Focused
two-guest tests isolate particular transitions; three-guest fixtures exercise
membership, placement, and failover; generic tests and safety models cover
larger or exhaustively explored configurations.

The next decisive evidence concerns resource-aware and rolling placement,
state transfer, production trust lifecycle, failure recovery, administration,
and sustained operations. Observing those mechanisms will show which
complexity genuinely disappears and which merely moves.

The intended endpoint is a cluster operating system with explicit
reconciliation and scheduling, acting on native, signed,
capability-constrained execution objects. At that boundary, the cluster stops
looking like a collection of computers on which software is installed and
starts looking like the computer itself.

\endinput
